\documentclass[11pt]{article}

\usepackage[a4paper,margin=2.7cm]{geometry}
\usepackage{amsmath,amssymb,amsthm,mathtools}
\usepackage{booktabs}
\usepackage{enumitem}

\newtheorem{definition}{Definition}
\newtheorem{assumption}{Assumption}
\newtheorem{theorem}{Theorem}
\newtheorem{proposition}{Proposition}
\newtheorem{remark}{Remark}
\newtheorem{example}{Example}

\newcommand{\Pcal}{\mathcal P}
\newcommand{\Mcal}{\mathcal M}
\newcommand{\E}{\mathbb E}
\newcommand{\R}{\mathbb R}
\newcommand{\W}{W_q}
\newcommand{\muhat}{\widehat\mu_n}
\newcommand{\norm}[1]{\left\lVert #1\right\rVert}

\title{Approximation of Regular Functionals of Probability Measures
by Interaction Polynomials}
\author{}
\date{}

\begin{document}

\maketitle

\section{Functionals of measures and Wasserstein polynomials}

Let \(K\subset\R^d\) be compact and let
\[
\Pcal(K)=\{\mu:\mu \text{ is a Borel probability measure on }K\}.
\]
Fix \(q\geq1\). We equip \(\Pcal(K)\) with the \(q\)-Wasserstein distance
\[
W_q(\mu,\nu)
=
\left(
\inf_{\pi\in\Pi(\mu,\nu)}
\int_{K\times K}|x-y|^q\,d\pi(x,y)
\right)^{1/q}.
\]

We study the approximation of a functional
\[
f:\Pcal(K)\longrightarrow\R.
\]
Unless stated otherwise, \(f\) is Borel measurable and bounded. For a
bounded functional \(g\), we use
\[
\norm{g}_\infty
=
\sup_{\mu\in\Pcal(K)}|g(\mu)|.
\]

\subsection{Interaction polynomials}

For \(m\geq 1\) and a bounded measurable kernel
\[
\Phi_m:K^m\longrightarrow\R,
\]
define
\[
P_{\Phi_m}(\mu)
=
\int_{K^m}
\Phi_m(x_1,\ldots,x_m)\,
d\mu(x_1)\cdots d\mu(x_m).
\]

Such a functional is called an interaction polynomial, or a Wasserstein
polynomial, of order at most \(m\).

Since \(\mu^{\otimes m}\) is invariant under permutations, the kernel may
always be symmetrized:
\[
\Phi_m^{\mathrm{sym}}(x_1,\ldots,x_m)
=
\frac1{m!}\sum_{\sigma\in\mathfrak S_m}
\Phi_m(x_{\sigma(1)},\ldots,x_{\sigma(m)}).
\]

The order \(m\) is the degree in the measure variable. It should not be
confused with the ordinary polynomial degree of \(\Phi_m\) in the spatial
variables \(x_1,\ldots,x_m\).

\subsection{Empirical-measure polynomial}

Given \(x=(x_1,\ldots,x_n)\in K^n\), set
\[
\mu_x^n=\frac1n\sum_{i=1}^n\delta_{x_i}.
\]
The canonical particle kernel associated with \(f\) is
\[
\Phi_n^f(x_1,\ldots,x_n)
=
f\left(\mu_x^n\right).
\]
The corresponding particle polynomial is
\[
B_nf(\mu)
=
\int_{K^n}
f\left(\frac1n\sum_{i=1}^n\delta_{x_i}\right)
d\mu^{\otimes n}(x_1,\ldots,x_n).
\]

Equivalently, if \(X_1,\ldots,X_n\) are independent with law \(\mu\), then
\[
B_nf(\mu)=\E f(\widehat\mu_n),
\qquad
\widehat\mu_n=\frac1n\sum_{i=1}^n\delta_{X_i}.
\]

Thus \(B_nf\) is a Wasserstein polynomial of order at most \(n\).

This operator has two complementary classical interpretations. If \(K\) is
finite, it is the multivariate Bernstein operator on the probability
simplex. For general compact \(K\), it is the canonical
Bernstein--Schnabl operator on the Bauer simplex \(\Pcal(K)\): one samples
the extreme points \(\delta_{X_i}\) and evaluates \(f\) at their barycentric
average. Probability-measure versions of this construction go back at
least to Schnabl \cite{schnabl}; the uniform convergence of the empirical
formulation for continuous \(f\) is also recorded in \cite{babu}.

In statistics, \(f(\widehat\mu_n)\) is the plug-in functional of the
empirical measure. The weak expansion of its expectation and the
higher-order particle correction used in
Section~\ref{sec:weak-correction} were developed systematically by
Chassagneux, Szpruch, and Tse
\cite{chassagneux-szpruch-tse}. Their analysis uses high-order linear
functional derivatives. Below we combine this weak-expansion mechanism
with a separate Wasserstein control of the Taylor remainder.

\section{Wasserstein--Hölder and Taylor regularity}

\subsection{Hölder regularity of order \(0<\alpha\leq 1\)}
\label{sec:metric-holder}

\begin{definition}[Metric Wasserstein--Hölder regularity]
For \(0<\alpha\leq1\), the functional \(f\) is
\(\alpha\)-Hölder with respect to \(W_q\) if
\[
|f(\mu)-f(\nu)|
\leq L W_q(\mu,\nu)^\alpha,
\qquad
\mu,\nu\in\Pcal(K).
\]
\end{definition}

\begin{proposition}[Hölder stability of the empirical operator]
\label{prop:holder-stability}
If \(f\) is \(\alpha\)-Hölder with constant \(L\), then, for every
\(n\geq1\) and \(\mu\in\Pcal(K)\),
\[
|B_nf(\mu)-f(\mu)|
\leq
L\E W_q(\widehat\mu_n,\mu)^\alpha.
\]
\end{proposition}

\begin{proof}
Using \(B_nf(\mu)=\E f(\widehat\mu_n)\), we obtain
\[
\begin{aligned}
|B_nf(\mu)-f(\mu)|
&=
\left|\E\bigl[f(\widehat\mu_n)-f(\mu)\bigr]\right|\\
&\leq
\E\left|f(\widehat\mu_n)-f(\mu)\right|\\
&\leq
L\E W_q(\widehat\mu_n,\mu)^\alpha.
\end{aligned}
\]
\end{proof}

A literal metric Hölder condition with exponent larger than one is too
restrictive on a geodesic metric space: subdividing a constant-speed
geodesic into \(N\) pieces shows that the total oscillation is bounded by
\(L N^{1-\alpha}\) times the \(\alpha\)-th power of its length, which
tends to zero. In particular, this applies to \((\Pcal(K),W_q)\) when
the underlying space is geodesic, for example when \(K\) is convex.
For a general compact \(K\), the Wasserstein space need not be geodesic;
in either case, higher differentiability is expressed by subtracting a
Taylor polynomial rather than by imposing a metric exponent above one.

\subsection{Wasserstein Taylor regularity of arbitrary order}
\label{sec:wasserstein-taylor}

For \(\alpha>1\), write
\[
\alpha=r+\theta,
\qquad
r=\lceil\alpha\rceil-1\in\mathbb N,
\qquad
0<\theta\leq 1.
\]
Thus an integer \(\alpha=m\) denotes the Lipschitz endpoint
\(C^{m-1,1}\).

Let \(\Mcal_0(K)\) be the vector space of finite signed Borel measures
on \(K\) with total mass zero. For \(\mu,\nu\in\Pcal(K)\), write
\[
\mu_t=(1-t)\mu+t\nu,
\qquad 0\leq t\leq1.
\]

\begin{definition}[Flat/vertical linear functional derivatives]
Set \(D^0f=f\). We say that \(f\) has linear functional derivatives up
to order \(r\) if there are symmetric \(j\)-linear forms
\[
D^jf(\mu):\Mcal_0(K)^j\longrightarrow\R,
\qquad 1\leq j\leq r,
\]
such that, for every mixture segment \(\mu_t\) and every fixed
\(\eta_1,\ldots,\eta_{j-1}\in\Mcal_0(K)\),
\[
\frac{d}{dt}
D^{j-1}f(\mu_t)[\eta_1,\ldots,\eta_{j-1}]
=
D^jf(\mu_t)
[\nu-\mu,\eta_1,\ldots,\eta_{j-1}].
\]
For \(j=1\), the empty argument list on \(D^0f=f\) is omitted.
Derivatives at \(t=0,1\) are understood one-sided. Thus
\(D^jf(\mu)[\eta_1,\ldots,\eta_j]\) is the mixed directional
derivative whenever the perturbations are admissible; the multilinear
form is assumed to extend to all of \(\Mcal_0(K)^j\).
Whenever continuity or an operator norm is invoked, the norm on
\(\Mcal_0(K)\) will be stated explicitly.
\end{definition}

These derivatives act along affine mixture paths: they change the mass
assigned to fixed points rather than transporting the points. They are
therefore iterated \emph{flat} or \emph{vertical} derivatives. More
precisely, at first order the definition above is the Gâteaux-like flat
derivative on \(\Pcal(K)\). A stronger Fréchet vertical derivative also
requires a uniform first-order remainder in a specified norm, such as
total variation; that stronger property is not implicit in the
definition above. We follow the terminology and distinctions summarized
in \cite{bertucci}.

When the derivatives admit kernels, we write
\[
D^jf(\mu)[\eta_1,\ldots,\eta_j]
=
\int_{K^j}
\frac{\delta^j f}{\delta\mu^j}
(\mu,x_1,\ldots,x_j)
\,d\eta_1(x_1)\cdots d\eta_j(x_j).
\]
The kernels are also called linear functional or von Mises derivatives;
this statistical functional calculus goes back to
\cite{von-mises} and is treated systematically in \cite{serfling}.
Since every \(\eta_i\) has zero mass, the kernel is not unique. We choose
the centred representative
\[
\left(\frac{\delta^j f}{\delta\mu^j}\right)_{\!c}
(\mu,x_1,\ldots,x_j)
=
D^jf(\mu)
[\delta_{x_1}-\mu,\ldots,\delta_{x_j}-\mu].
\]
It represents the same multilinear form and satisfies
\[
\int_K
\left(\frac{\delta^j f}{\delta\mu^j}\right)_{\!c}
(\mu,x_1,\ldots,x_j)\,d\mu(x_i)
=0
\]
for each \(i\), with the remaining variables held fixed.
We omit the subscript \(c\) for this representative below.

There is an important distinction from the horizontal, or Lions,
derivative. If \(K\subset\R^d\) and the first vertical kernel is
sufficiently regular in \(x\), then
\[
D_\mu^{\mathrm H}f(\mu,x)
=
\nabla_x\frac{\delta f}{\delta\mu}(\mu,x).
\]
The vector-valued derivative on the left tests transport perturbations
and, under the standard lift to random variables, agrees with the Lions
derivative. Thus the scalar vertical derivative and the Lions derivative
are not two names for the same object: the latter is obtained from the
former by a spatial gradient when the required regularity is available.
See \cite{bertucci,carmona-delarue}. All Taylor and particle expansions
in this note use the vertical derivatives.

The Taylor polynomial of order \(r\) around \(\mu\) is
\[
T_\mu^r f(\nu)
=
f(\mu)
+
\sum_{j=1}^r
\frac1{j!}
D^jf(\mu)
\left[(\nu-\mu)^{\otimes j}\right].
\]

\begin{definition}[Wasserstein--Taylor class]
We say that \(f\in C_W^\alpha\) if its Taylor polynomial of order \(r\)
satisfies
\[
\left|
f(\nu)-T_\mu^r f(\nu)
\right|
\leq
L W_q(\mu,\nu)^\alpha
\]
uniformly over \(\mu,\nu\in\Pcal(K)\).
\end{definition}

The remainder estimate is the definition; an operator norm must be
specified before stating a derivative-level criterion. On
\(\Mcal_0(K)\), define the Kantorovich--Rubinstein norm
\[
\norm{\eta}_{\mathrm{KR}}
=
\sup_{\operatorname{Lip}(\varphi)\leq1}
\left|\int_K\varphi\,d\eta\right|.
\]
In particular,
\[
\norm{\nu-\mu}_{\mathrm{KR}}
=W_1(\mu,\nu)
\leq W_q(\mu,\nu).
\]
For an \(r\)-linear form \(A\), use the induced operator norm
\[
\norm{A}_{\mathrm{op},\mathrm{KR}}
=
\sup_{\norm{\eta_i}_{\mathrm{KR}}\leq1}
\left|A[\eta_1,\ldots,\eta_r]\right|.
\]
This is not the \(L^\infty(K^r)\) norm of a derivative kernel: a kernel
supremum norm controls total-variation increments and does not supply
the \(r\) powers of \(W_q\) needed below.

A sufficient condition for membership in \(C_W^\alpha\) is
\[
\norm{D^rf(\mu)-D^rf(\nu)}_{\mathrm{op},\mathrm{KR}}
\leq
L W_q(\mu,\nu)^\theta.
\]
Indeed, with \(\Delta=\nu-\mu\), the integral Taylor formula gives
\[
f(\nu)-T_\mu^rf(\nu)
=
\frac1{(r-1)!}
\int_0^1(1-t)^{r-1}
\bigl(D^rf(\mu_t)-D^rf(\mu)\bigr)
[\Delta^{\otimes r}]\,dt.
\]
Using the operator norm,
\[
\begin{aligned}
|f(\nu)-T_\mu^rf(\nu)|
&\leq
\frac{L}{(r-1)!}
\int_0^1
(1-t)^{r-1}
W_q(\mu_t,\mu)^\theta
\norm{\Delta}_{\mathrm{KR}}^r\,dt.
\end{aligned}
\]
The coupling that leaves a mass \(1-t\) fixed and transports the
remaining mass gives
\[
W_q(\mu_t,\mu)
\leq t^{1/q}W_q(\mu,\nu).
\]
Consequently,
\[
|f(\nu)-T_\mu^rf(\nu)|
\leq
C_{r,\theta,q}L
W_q(\mu,\nu)^\theta W_1(\mu,\nu)^r
\leq
C_{r,\theta,q}L W_q(\mu,\nu)^{r+\theta}.
\]
The converse need not hold without additional extension and
differentiability assumptions, which is why this criterion is not
stated as an equivalence.

This is the analogue on \(\Pcal(K)\) of the classical decomposition
\[
\alpha=r+\theta
\]
for functions in a Hölder space \(C^{r,\theta}\).

\begin{remark}[Relation with \(C^{1,\theta}\) terminology]
When \(r=1\), the fractional exponent in the derivative-level condition
is \(\theta\), and the total Taylor order is
\(\alpha=1+\theta\). Thus the classes denoted \(C^{1,\theta}\) in
\cite{bertucci} have total order \(1+\theta\) in the convention of this
note. They should not, however, be identified automatically with
\(C_W^{1+\theta}\): Bertucci's vertical class uses total variation in
the measure variable and the supremum norm of the flat kernel, while his
horizontal class controls the Lions derivative along couplings. Our
sufficient criterion instead combines a vertical derivative, the
Kantorovich--Rubinstein operator norm, and a \(W_q^\theta\) modulus.
\end{remark}

\subsection{Taylor regularity in a root-\(n\) tangent norm}
\label{sec:root-n-taylor}

The functional derivatives are the same linear derivatives as above.
What changes is the geometry used to control the remainder. We keep
\(r\) for the Taylor degree throughout, rather than introducing a
second symbol.

Fix an integer \(r\geq1\). Suppose that
\[
f(\nu)
=
f(\mu)
+
\sum_{j=1}^{r}
\frac1{j!}
D^jf(\mu)
\left[(\nu-\mu)^{\otimes j}\right]
+
R_{r+1}(\mu,\nu),
\]
and that centred derivative kernels can be chosen so that
\[
\sup_{\mu\in\Pcal(K)}
\norm{
\frac{\delta^j f}{\delta\mu^j}(\mu,\cdot)
}_{L^\infty(K^j)}
<\infty,
\qquad 1\leq j\leq r.
\]
This boundedness is used later to control the expected Taylor
contractions; the root-\(n\) remainder estimate itself uses only the
next two assumptions.

Let \(\norm{\cdot}_{\mathcal T}\) be a norm or seminorm on
\(\Mcal_0(K)\) for which the empirical fluctuation has root-\(n\)
moments:
\[
\sup_{\mu\in\Pcal(K)}
\E\norm{\widehat\mu_n-\mu}_{\mathcal T}^{r+1}
\leq
C_r n^{-(r+1)/2}.
\]
Assume that
\[
|R_{r+1}(\mu,\nu)|
\leq
C\norm{\nu-\mu}_{\mathcal T}^{r+1}.
\]
Then
\[
\sup_{\mu\in\Pcal(K)}
\E|R_{r+1}(\mu,\widehat\mu_n)|
\leq
CC_r n^{-(r+1)/2}.
\]

The comparison with a Wasserstein remainder requires an explicit
relation between the two geometries. In the typical situation
\[
\norm{\nu-\mu}_{\mathcal T}
\leq C_{\mathcal T}W_q(\mu,\nu),
\]
the tangent norm is weaker, so the upper bound
\(C\norm{\nu-\mu}_{\mathcal T}^{r+1}\) is smaller and therefore more
demanding than a \(W_q^{r+1}\) remainder bound. The tangent estimate
implies the Wasserstein one, but the converse generally fails. Along
empirical measures, the tangent hypothesis gives the dimension-free
scale \(n^{-(r+1)/2}\), whereas a bare Wasserstein-remainder hypothesis
gives \(\rho_{n,q,d}^{\,r+1}\), defined in
Section~\ref{sec:raw-particle} below.

For example, given finitely many bounded Lipschitz functions
\(\varphi_1,\ldots,\varphi_L\), the finite-moment seminorm
\[
\norm{\eta}_{\Phi}
=
\left(
\sum_{\ell=1}^L
\left|\int_K\varphi_\ell\,d\eta\right|^2
\right)^{1/2}
\]
has root-\(n\) empirical moments and satisfies
\(\norm{\nu-\mu}_{\Phi}\lesssim W_1(\mu,\nu)\leq W_q(\mu,\nu)\).
It may vanish for two distinct measures having the same selected
moments. Requiring the entire Taylor remainder to be controlled by this
much smaller quantity is therefore a genuine structural restriction on
\(f\), not merely a smoother choice of notation.

\section{Examples of functionals}

\subsection{Linear moments}

For a bounded function \(V:K\to\R\), let
\[
f(\mu)=\int_K V(x)\,d\mu(x).
\]
This is a polynomial of order one. Moreover,
\[
B_nf(\mu)=f(\mu)
\]
for every \(n\), so the approximation error is zero.

\subsection{Pair interactions}

Let \(W:K\times K\to\R\) be bounded and Borel measurable, and define
\[
f(\mu)
=
\iint_{K\times K}
W(x,y)\,d\mu(x)d\mu(y).
\]
This is a polynomial of order two. The particle approximation satisfies
the exact identity
\[
B_nf(\mu)-f(\mu)
=
\frac1n
\left[
\int_K W(x,x)\,d\mu(x)
-
\iint_{K\times K}W(x,y)\,d\mu(x)d\mu(y)
\right].
\]

The diagonal bias can be removed by the \(U\)-statistic
\[
\frac1{n(n-1)}
\sum_{i\neq j}W(X_i,X_j).
\]
This is the classical passage from a \(V\)-statistic to a
\(U\)-statistic; see \cite{hoeffding,serfling}.

\subsection{Smooth functions of finitely many moments}

Let \(\varphi_1,\ldots,\varphi_L:K\to\R\) be bounded and Lipschitz, and
let \(F:\R^L\to\R\). Consider
\[
f(\mu)
=
F\left(
\int_K\varphi_1\,d\mu,\ldots,
\int_K\varphi_L\,d\mu
\right).
\]

If \(F\in C^{r,\theta}\) on a neighborhood of the compact set of
attainable moment vectors \(m_\mu\), with bounded derivatives there,
then \(f\) has functional derivatives of the same order. For example,
\[
\frac{\delta f}{\delta\mu}(\mu,x)
=
\sum_{\ell=1}^L
\partial_\ell F(m_\mu)
\left(
\varphi_\ell(x)-\int_K\varphi_\ell\,d\mu
\right),
\]
where
\[
m_\mu
=
\left(
\int\varphi_1\,d\mu,\ldots,\int\varphi_L\,d\mu
\right).
\]

These are prototypical high-regularity functionals for which empirical
fluctuations are finite dimensional and therefore of order \(n^{-1/2}\),
independently of \(d\).

\subsection{Squared maximum mean discrepancy}

Let \(k\) be a bounded positive-definite kernel with reproducing kernel
Hilbert space \(H\), and define
\[
m_\mu=\int_K k(x,\cdot)\,d\mu(x).
\]
For a fixed \(\nu\), let
\[
f_\nu(\mu)
=
\operatorname{MMD}_k^2(\mu,\nu)
=
\norm{m_\mu-m_\nu}_H^2.
\]

Expanding the square gives
\[
f_\nu(\mu)
=
\iint k(x,y)\,d\mu(x)d\mu(y)
-
2\iint k(x,y)\,d\mu(x)d\nu(y)
+
C_\nu.
\]
Thus squared MMD is exactly a polynomial of order two.

Its particle bias is
\[
B_nf_\nu(\mu)-f_\nu(\mu)
=
\frac1n
\left[
\int k(x,x)\,d\mu(x)
-
\iint k(x,y)\,d\mu(x)d\mu(y)
\right].
\]

The corresponding \(U\)-statistic removes this bias exactly.

The unsquared MMD,
\[
\operatorname{MMD}_k(\mu,\nu)
=
\norm{m_\mu-m_\nu}_H,
\]
is not differentiable at \(\mu=\nu\). At the null,
\[
\E\operatorname{MMD}_k(\widehat\mu_n,\mu)
=
O(n^{-1/2}),
\]
whereas away from the null a Taylor expansion can give a weak bias of
order \(1/n\).

\subsection{Squared Wasserstein distance}

Let
\[
f_\nu(\mu)=W_2^2(\mu,\nu).
\]
If \(D=\operatorname{diam}(K)\), then
\[
\left|
W_2^2(\mu,\nu)-W_2^2(\mu',\nu)
\right|
\leq
2D W_2(\mu,\mu').
\]
Therefore
\[
|B_nf_\nu(\mu)-f_\nu(\mu)|
\leq
2D\E W_2(\widehat\mu_n,\mu).
\]

The functional \(W_2^2(\cdot,\nu)\) is not globally smooth in the
von Mises sense. Its differentiability depends on uniqueness and
regularity of the optimal transport map.

At \(\nu=\mu\),
\[
B_nf_\mu(\mu)-f_\mu(\mu)
=
\E W_2^2(\widehat\mu_n,\mu).
\]
Uniformly over compactly supported measures,
\[
\E W_2^2(\widehat\mu_n,\mu)
\lesssim
\begin{cases}
n^{-1/2},&d<4,\\[1mm]
n^{-1/2}\log(1+n),&d=4,\\[1mm]
n^{-2/d},&d>4.
\end{cases}
\]
These uniform compact-support bounds follow, for example, from the
empirical Wasserstein estimates in
\cite{dereich-scheutzow-schottstedt,fournier-guillin}.

There are special linear cases. For example,
\[
W_2^2(\mu,\delta_a)
=
\int_K|x-a|^2\,d\mu(x),
\]
and consequently the particle approximation is exact.

\subsection{Relative entropy}

This example is a warning outside the standing bounded real-valued
setting: without restrictions or regularization, relative entropy is
extended valued.

For
\[
f_\nu(\mu)=\operatorname{KL}(\mu\Vert\nu),
\]
the direct empirical kernel is generally not finite. If \(\nu\) is
nonatomic and has a density, then
\[
\widehat\mu_n\not\ll\nu
\]
and hence
\[
\operatorname{KL}(\widehat\mu_n\Vert\nu)=+\infty
\]
almost surely.

On a finite state space with \(S\) states,
\[
f_{p^\star}(p)
=
\sum_{j=1}^S p_j\log\frac{p_j}{p_j^\star}
\]
is smooth in the interior of the simplex. If all probabilities are
bounded away from zero, then
\[
\E f_{p^\star}(\widehat p_n)-f_{p^\star}(p)
=
\frac{S-1}{2n}+O(n^{-2}).
\]

For continuous measures, one may instead regularize:
\[
f_{\nu,\varepsilon}(\mu)
=
\operatorname{KL}(K_\varepsilon\mu\Vert\nu),
\]
where \(K_\varepsilon\mu\) is a smoothed density. For fixed
\(\varepsilon>0\), this functional can possess bounded functional
derivatives and a \(1/n\) weak expansion. The constants generally grow
as a power of \(\varepsilon^{-d}\).

\section{Rate of the particle interaction polynomial}

\subsection{Empirical Wasserstein scale and the raw particle error}
\label{sec:raw-particle}

Define the characteristic empirical Wasserstein scale
\[
\rho_{n,q,d}
=
\begin{cases}
n^{-1/(2q)},&d<2q,\\[1mm]
n^{-1/(2q)}(\log(1+n))^{1/q},&d=2q,\\[1mm]
n^{-1/d},&d>2q.
\end{cases}
\]

For every fixed \(a>0\), compact support gives the uniform moment
estimate
\[
\sup_{\mu\in\Pcal(K)}
\left(
\E W_q(\widehat\mu_n,\mu)^a
\right)^{1/a}
\leq
C_{a,q,d,K}\rho_{n,q,d};
\]
this is a uniform ambient-dimensional upper bound
\cite{dereich-scheutzow-schottstedt,fournier-guillin}. It need not be
sharp when \(K\) is finite or has smaller intrinsic dimension.

The smooth weak-bias component of the result below is the static
i.i.d.\ empirical-measure expansion developed in
\cite{chassagneux-szpruch-tse}. The proposition combines that
cancellation mechanism with the Wasserstein remainder, thereby also
covering fractional and dimension-dependent regularity regimes.

\begin{proposition}[Raw particle error across regularity regimes]
Let \(\alpha>0\). Assume the metric Hölder condition of
Section~\ref{sec:metric-holder}
when \(0<\alpha\leq1\). When \(\alpha>1\), assume
\(f\in C_W^\alpha\), with \(r=\lceil\alpha\rceil-1\), and assume that
each \(D^jf\), \(1\leq j\leq r\), admits a centred kernel that is
uniformly bounded over \(\mu\in\Pcal(K)\). Then
\[
\boxed{
\norm{B_nf-f}_\infty
\lesssim
\begin{cases}
\rho_{n,q,d}^{\,\alpha},
&0<\alpha\leq2,\\[1mm]
n^{-1}+\rho_{n,q,d}^{\,\alpha},
&\alpha>2.
\end{cases}
}
\]
\end{proposition}

\begin{proof}
For \(0<\alpha\leq1\), Proposition~\ref{prop:holder-stability} and the
empirical moment estimate give the result. Let \(\alpha>1\), put
\[
\Delta_n=\widehat\mu_n-\mu,
\qquad
Z_i=\delta_{X_i}-\mu,
\]
and use the Taylor expansion
\[
f(\widehat\mu_n)-f(\mu)
=
\sum_{j=1}^r
\frac1{j!}D^jf(\mu)[\Delta_n^{\otimes j}]
+
R_\alpha(\mu,\widehat\mu_n).
\]
The first-order term is centred:
\[
\E Df(\mu)[\Delta_n]=0.
\]
For \(j\geq2\), multilinearity gives
\[
\E D^jf(\mu)[\Delta_n^{\otimes j}]
=
\frac1{n^j}
\sum_{i_1,\ldots,i_j=1}^n
\E D^jf(\mu)[Z_{i_1},\ldots,Z_{i_j}].
\]
A summand vanishes if one particle index occurs exactly once, by
independence and \(\E Z_i=0\). A surviving tuple therefore contains at
most \(\lfloor j/2\rfloor\) distinct indices. Uniform boundedness of
the derivative kernel then yields
\[
\left|
\E D^jf(\mu)[\Delta_n^{\otimes j}]
\right|
\lesssim
n^{-j+\lfloor j/2\rfloor}
=
n^{-\lceil j/2\rceil}.
\]
Finally,
\[
\E|R_\alpha(\mu,\widehat\mu_n)|
\lesssim
\rho_{n,q,d}^{\,\alpha}.
\]
For \(1<\alpha\leq2\), one has \(r=1\), so only the centred linear term
appears. For \(\alpha>2\), all terms of order \(j\geq2\) are
\(O(n^{-1})\).
\end{proof}

In the high-dimensional regime \(d>2q\), the proposition becomes
\[
\boxed{
\norm{B_nf-f}_\infty
\lesssim
\begin{cases}
n^{-\alpha/d},&0<\alpha\leq2,\\[1mm]
n^{-\min\{1,\alpha/d\}},&\alpha>2.
\end{cases}
}
\]
Thus Section~\ref{sec:raw-particle} is not restricted to
\(\alpha\leq1\): the first line
uses metric Hölder regularity for \(0<\alpha\leq1\) and the
first-order Taylor remainder for \(1<\alpha\leq2\). The second line
uses a genuine second derivative and higher-order Taylor regularity.

\subsection{The quadratic bias and the \(1/n\) threshold}

When \(D^2f\) exists, the quadratic contraction can be written exactly
as
\[
\frac12\E D^2f(\mu)[\Delta_n^{\otimes2}]
=
\frac{A_1f(\mu)}{n},
\]
where
\[
A_1f(\mu)
=
\frac12
\int_K
D^2f(\mu)
[\delta_x-\mu,\delta_x-\mu]\,d\mu(x).
\]
Under sufficiently high linear, or von Mises, Taylor regularity with a
remainder smaller than the retained terms, the remaining contractions
have an expansion in integer powers of \(1/n\), beginning with
\[
B_nf(\mu)-f(\mu)
=
\frac{A_1f(\mu)}{n}
+
\frac{A_2f(\mu)}{n^2}
+\cdots.
\]
This is the classical von Mises expansion of a plug-in functional
\cite{von-mises,serfling}. An all-order treatment of this static
empirical-measure expansion under high-order linear functional
derivatives, together with particle applications, is given in
\cite{chassagneux-szpruch-tse}.

This explains the saturation statement precisely. If the terms beyond
\(A_1f/n\) are \(o(n^{-1})\) and \(A_1f\) is not identically zero, the
uncorrected operator is generically of exact order \(1/n\); additional
derivatives do not cancel \(A_1f\). For example, for
\[
f(\mu)
=
\left(\int_K\varphi\,d\mu\right)^2,
\]
one has the exact identity
\[
B_nf(\mu)-f(\mu)
=
\frac{\operatorname{Var}_\mu(\varphi)}{n},
\]
although \(f\) is already a polynomial and has derivatives of every
order. Linear functionals and other special cases may have
\(A_1f=0\), so \(1/n\) is not a universal lower bound.

Under Wasserstein--Taylor regularity alone, the remainder
\(\rho_{n,q,d}^{\,\alpha}\) may be larger than \(1/n\). Thus
\(\alpha=2\) by itself does not imply a \(1/n\) rate. In the regime
\(d>2q\), the Wasserstein remainder is \(O(n^{-1})\) only when
\(\alpha\geq d\). For \(d<2q\), the corresponding condition is
\(\alpha\geq2q\); at the critical dimension \(d=2q\), equality still
carries a logarithmic loss. By contrast, the root-\(n\) tangent
regularity of Section~\ref{sec:root-n-taylor} gives a dimension-free
\(O(n^{-1})\) raw bound already from a quadratic tangent remainder.

\subsection{Higher-order corrected construction}
\label{sec:weak-correction}

Let \(J\geq2\) denote the desired \emph{weak correction order}; it is
not the smoothness index \(\alpha\). Suppose that
\(r=\lceil\alpha\rceil-1\geq2J-2\), the derivative kernels through
order \(r\) are uniformly bounded, and the \(C_W^\alpha\) remainder
bound holds. The same index-partition argument as above groups the
Taylor contractions into
\[
B_mf
=
f
+
\sum_{\ell=1}^{J-1}\frac{A_\ell f}{m^\ell}
+
\mathcal E_{m,J},
\]
with
\[
\norm{\mathcal E_{m,J}}_\infty
\leq
C\left(m^{-J}+\rho_{m,q,d}^{\,\alpha}\right).
\]
Indeed, the coefficient \(A_\ell f\) involves derivatives only through
order \(2\ell\); terms whose first possible weak order is at least
\(J\) enter \(\mathcal E_{m,J}\). The strict condition
\(\alpha>2J-2\) is an equivalent convenient way to write
\(r\geq2J-2\) under the convention of
Section~\ref{sec:wasserstein-taylor}.

Choose distinct positive integers
\[
\lambda_1,\ldots,\lambda_J
\]
and coefficients \(a_1,\ldots,a_J\) satisfying
\[
\sum_{j=1}^J a_j=1,
\qquad
\sum_{j=1}^J a_j\lambda_j^{-\ell}=0,
\quad
1\leq\ell\leq J-1.
\]

This Vandermonde system has a unique solution. Define the
Richardson--Romberg particle polynomial
\[
\mathcal B_{n,J}f
=
\sum_{j=1}^J a_j B_{\lambda_j n}f.
\]
This is Richardson extrapolation \cite{richardson} applied to the weak
particle expansion; its use in this measure-functional setting follows
\cite{chassagneux-szpruch-tse}.
It is a signed interaction polynomial of order at most
\[
\max_j\lambda_j n.
\]
For instance, \(J=2\), \((\lambda_1,\lambda_2)=(1,2)\), gives
\[
\mathcal B_{n,2}f=2B_{2n}f-B_nf,
\]
which cancels the \(A_1f/n\) term.

In general, the first \(J-1\) weak bias terms cancel, and
\[
\boxed{
\norm{\mathcal B_{n,J}f-f}_\infty
\lesssim
n^{-J}+\rho_{n,q,d}^{\,\alpha}.
}
\]

In the high-dimensional regime \(d>2q\),
\[
\boxed{
\norm{\mathcal B_{n,J}f-f}_\infty
\lesssim
n^{-\min\{J,\alpha/d\}}.
}
\]

The minimum has three distinct meanings:
\begin{itemize}
\item If \(\alpha/d\leq1\), the Wasserstein remainder term already
controls the general raw-operator bound, so correction does not improve
the guaranteed exponent.
\item If \(1<\alpha/d<J\), correction removes the \(1/n\) saturation
and the bound reaches the full \(n^{-\alpha/d}\) Wasserstein remainder
rate.
\item If \(\alpha/d\geq J\), the corrected weak-bias term \(n^{-J}\)
controls the displayed bound; increasing \(J\) improves the guarantee
if the required derivatives are available.
\end{itemize}
Thus, when \(J\geq\alpha/d\), the minimum is indeed \(\alpha/d\), but
this means that the construction has successfully removed the weak
bias term from this upper bound and reached the general Wasserstein
remainder ceiling. It does not mean that the correction failed to use
smoothness. For example, if \(\alpha/d=2.4\), the guaranteed raw
exponent is \(1\), correction order \(J=2\) gives exponent \(2\), and
\(J=3\) reaches exponent \(2.4\). Special functionals with vanishing
contractions or a sharper remainder can, of course, converge faster
than these general bounds.

\begin{remark}[Relation with the Remez algorithm]
The correction above is neither the Remez algorithm nor a direct extension
of it. Classical Remez approximation fixes a polynomial degree \(p\) and
computes a minimax approximant
\[
q_p^\star
\in
\operatorname*{arg\,min}_{q\in\Pi_p}
\norm{g-q}_{L^\infty([a,b])},
\]
where \(\Pi_p\) is the space of univariate polynomials of degree at most
\(p\). Its exchange iteration enforces the Chebyshev equioscillation
condition at \(p+2\) adaptively chosen extremal points
\cite{remez}. By contrast, the coefficients \(a_j\) in
\(\mathcal B_{n,J}\) are fixed by cancelling the first \(J-1\) powers in
the weak particle-bias expansion. There is no exchange of extremal
measures, no equioscillation condition, and no claim that
\(\mathcal B_{n,J}f\) is a best uniform approximation among interaction
polynomials of a prescribed order. The limited analogy is that both
methods determine signed coefficients through a structured linear system
in order to reduce a uniform approximation error.

A genuine Remez-type problem on \(\Pcal(K)\) would first choose a
finite-dimensional space \(\mathcal V_N\) of interaction or moment
polynomials and solve
\[
P_N^\star
\in
\operatorname*{arg\,min}_{P\in\mathcal V_N}
\norm{f-P}_\infty.
\]
That is a different, generally infinite-dimensional minimax problem unless
\(\mathcal V_N\) is fixed explicitly. Remez is more directly relevant to
the second approximation level of Section~\ref{sec:two-level}: for a
fixed particle order \(n\), it may be used to compute a minimax spatial
kernel \(Q_{n,p}\) when that kernel depends on a single scalar variable,
possibly after a structural reduction. For a generic kernel on
\(K^n\subset\R^{nd}\), the classical one-dimensional equioscillation
theorem does not apply directly. Thus the
particle/Richardson stage exploits smoothness in the measure variable,
whereas a Remez step could optimize the separate spatial-kernel
approximation.
\end{remark}

\subsection{Dimension-free weak Taylor rate}

Suppose instead that the degree-\(r\) Taylor expansion satisfies the
root-\(n\) tangent assumptions of Section~\ref{sec:root-n-taylor}. Put
\[
\sigma_r=\frac{r+1}{2},
\qquad
J_r=\lceil \sigma_r\rceil.
\]
Each expected Taylor contraction is a finite sum of integer powers of
\(1/n\), and its smallest possible power is
\(\lceil j/2\rceil\) at derivative order \(j\). Consequently,
\[
B_mf
=
f
+
\sum_{1\leq\ell<\sigma_r}\frac{A_\ell f}{m^\ell}
+
\mathcal R_{m,r},
\qquad
\norm{\mathcal R_{m,r}}_\infty
\lesssim
m^{-\sigma_r}.
\]
Using \(J_r\) Richardson levels and cancelling the powers
\(\ell=1,\ldots,J_r-1\) therefore gives, for the previously fixed
levels \(\lambda_j\), the corrected operator
\(\mathcal B_{n,J_r}\), with the convention
\(\mathcal B_{n,1}=B_n\), and the dimension-free bound
\[
\boxed{
\norm{\mathcal B_{n,J_r}f-f}_\infty
\lesssim
n^{-(r+1)/2}.
}
\]
There is no need to replace this by a floor exponent: the possible
half-integer rate comes from the explicitly assumed
\((r+1)\)-st moment of the tangent remainder. Roughly two additional
bounded functional derivatives are needed for each additional power of
\(1/n\). Without correction, the generic rate is \(1/n\): for
\(r\geq2\) the quadratic contraction remains, while for \(r=1\) the
assumed quadratic tangent remainder itself is \(O(n^{-1})\).

\section{Two-level approximation of the interaction kernel}
\label{sec:two-level}

Let \(P_nf\) denote a chosen first-level interaction polynomial of
order at most \(n\). For \(B_nf\), this is its natural particle order.
For a corrected operator \(\mathcal B_{m,J}\), set
\[
\Lambda_J=\max_{1\leq j\leq J}\lambda_j,
\qquad
n=\Lambda_Jm,
\]
and lift every lower-order term by adding dummy variables; the result
is then \(P_nf=\mathcal B_{m,J}f\). Since \(J\) and the levels
\(\lambda_j\) are fixed, \(n\asymp m\), so this reindexing does not
change any algebraic rate. Write the resulting kernel as
\[
P_n f(\mu)
=
\int_{K^n}
\Phi_n(x_1,\ldots,x_n)\,d\mu^{\otimes n}.
\]

Although each particle belongs to \(\R^d\), the complete kernel is a
function on
\[
K^n\subset\R^{nd}.
\]
Approximating a generic kernel is therefore an \(nd\)-dimensional
approximation problem.

\subsection{Ordinary polynomial approximation of the kernel}

Let \(Q_{n,p}\) be an ordinary polynomial in the \(nd\) scalar
variables, using a specified tensor-degree or total-degree convention
with degree parameter \(p\). Define
\[
P_{n,p}f(\mu)
=
\int_{K^n}
Q_{n,p}(x_1,\ldots,x_n)\,d\mu^{\otimes n}.
\]

Assume that \(\Phi_n\) has spatial smoothness \(s>0\) and that, for
some \(\gamma\geq0\),
\[
\boxed{
\norm{\Phi_n-Q_{n,p}}_\infty
\leq
C n^\gamma p^{-s}.
}
\]
The factor \(n^\gamma\) allows the spatial smoothness norm of the kernel
to grow with the number of particles.

Measure-variable regularity and spatial regularity are distinct:
bounded functional derivatives of \(f\) do not automatically imply
uniform \(C^s\)-regularity of \(\Phi_n\) as a function on \(K^n\).

Suppose the first approximation level satisfies
\[
\norm{f-P_nf}_\infty\leq Cn^{-a}.
\]
Then
\[
\begin{aligned}
\norm{f-P_{n,p}f}_\infty
&\leq
\norm{f-P_nf}_\infty
+
\norm{\Phi_n-Q_{n,p}}_\infty\\
&\leq
C_0n^{-a}+C_1n^\gamma p^{-s}.
\end{aligned}
\]

Consequently,
\[
\boxed{
\norm{f-P_{n,p}f}_\infty
\leq
C_0n^{-a}+C_1n^\gamma p^{-s}.
}
\]

The relevant first-level exponent is, for example,
\[
a=
\begin{cases}
\alpha/d,
&
\text{for the unsaturated high-dimensional rate, }0<\alpha\leq2,\\[1mm]
\min\{1,\alpha/d\},
&
\text{for the uncorrected high-dimensional rate when \(\alpha>2\)},\\[1mm]
\min\{J,\alpha/d\},
&
\text{for weak correction order }J\text{, assuming }\alpha>2J-2,\\[1mm]
1,
&
\text{for an uncorrected root-\(n\) tangent remainder},\\[1mm]
(r+1)/2,
&
\text{for a corrected root-\(n\) Taylor expansion of degree \(r\)}.
\end{cases}
\]
In the first row, regularity means metric Hölder regularity when
\(0<\alpha\leq1\) and Wasserstein--Taylor regularity when
\(1<\alpha\leq2\).

\subsection{Optimal choice of \(p\) for a fixed \(n\)}

Balancing the two errors,
\[
n^{-a}\asymp n^\gamma p^{-s},
\]
gives
\[
\boxed{
p\asymp n^{(a+\gamma)/s}.
}
\]

With this choice,
\[
\boxed{
\norm{f-P_{n,p}f}_\infty
\lesssim n^{-a}.
}
\]

In the stable case \(\gamma=0\),
\[
p\asymp n^{a/s}.
\]

For an error tolerance \(\varepsilon\),
\[
n\asymp\varepsilon^{-1/a},
\qquad
p\asymp
\varepsilon^{-(a+\gamma)/(as)}.
\]
When \(\gamma=0\),
\[
p\asymp\varepsilon^{-1/s}.
\]

For example, suppose the first-level construction reaches the
high-dimensional Wasserstein remainder exponent \(\alpha/d\), as it
does for \(0<\alpha\leq2\), or after correction when
\(J\geq\alpha/d\) under the hypotheses of
Section~\ref{sec:weak-correction}. In the stable case \(\gamma=0\), if
\[
a=\frac{\alpha}{d}
\qquad\text{and}\qquad
s=\alpha,
\]
then
\[
\boxed{
p\asymp n^{1/d},
\qquad
\norm{f-P_{n,p}f}_\infty
\lesssim n^{-\alpha/d}.
}
\]

Likewise, if \(\gamma=0\) and the corrected tangent Taylor rate is
\[
a=\frac{r+1}{2}
\]
and the spatial smoothness is \(s=r+1\), then
\[
\boxed{
p\asymp n^{1/2},
\qquad
\norm{f-P_{n,p}f}_\infty
\lesssim n^{-(r+1)/2}.
}
\]

\subsection{Representation in terms of ordinary moments}

Write the polynomial kernel as
\[
Q_{n,p}(x_1,\ldots,x_n)
=
\sum_{\ell=1}^M
c_\ell
\prod_{i=1}^n x_i^{\,a_{\ell,i}},
\]
where
\[
a_{\ell,i}\in\mathbb N^d,
\qquad
x_i^{\,a_{\ell,i}}
=
\prod_{j=1}^d x_{i,j}^{a_{\ell,i,j}}.
\]

Integration against \(\mu^{\otimes n}\) gives
\[
P_{n,p}f(\mu)
=
\sum_{\ell=1}^M
c_\ell
\prod_{i=1}^n
m_{a_{\ell,i}}(\mu),
\]
where
\[
m_a(\mu)=\int_K x^a\,d\mu(x)
\]
is an ordinary moment of \(\mu\).

Thus the two-level construction eventually produces an ordinary
polynomial in finitely many moments.

\subsection{Number of coefficients}

A full tensor polynomial with degree at most \(p\) in each of the
\(nd\) scalar variables has
\[
\boxed{
M_{\mathrm{tensor}}(n,p)
=
(p+1)^{nd}
}
\]
coefficients.

A polynomial of total degree at most \(p\) has
\[
\boxed{
M_{\mathrm{total}}(n,p)
=
\binom{nd+p}{p}.
}
\]

For a generic \(s\)-smooth function of \(nd\) variables, the classical
fixed-dimensional approximation relation can be expressed as
\[
\text{kernel error}
\lesssim
M^{-s/(nd)}.
\]
Using this estimate while \(n\) varies is an additional uniformity
assumption: all constants depending on the growing dimension \(nd\)
must be controlled by the displayed factor \(n^\gamma\). Such control
does not follow from the fixed-dimensional theorem alone.

Let \(P_{n,M}f\) denote the resulting two-level approximation when its
ordinary polynomial kernel has \(M\) coefficients. A parameter-count
version of the estimate is
\[
\boxed{
\norm{f-P_{n,M}f}_\infty
\lesssim
n^{-a}
+
n^\gamma M^{-s/(nd)}.
}
\]

This formula displays the principal difficulty: the exponent
\(s/(nd)\) deteriorates as the particle order \(n\) increases.

\subsection{Optimal rate as a function of the total parameter count}

Balancing
\[
n^{-a}
\asymp
n^\gamma M^{-s/(nd)}
\]
gives
\[
(a+\gamma)\log n
\asymp
\frac{s}{nd}\log M.
\]
Equivalently,
\[
n\log n
\asymp
\frac{s}{d(a+\gamma)}\log M.
\]

Hence
\[
n
\asymp
\frac{\log M}{\log\log M},
\]
up to constants depending on \(a,s,d,\gamma\). Substitution into the
first-level error gives
\[
\boxed{
\norm{f-P_{n,M}f}_\infty
\lesssim
\left(
\frac{\log M}{\log\log M}
\right)^{-a}.
}
\]

If the first-level construction reaches the high-dimensional
Wasserstein remainder rate---for instance, when
\(0<\alpha\leq2\), or after correction with \(J\geq\alpha/d\) under
the hypotheses of Section~\ref{sec:weak-correction}---then
\[
a=\frac{\alpha}{d},
\]
this becomes
\[
\boxed{
\norm{f-P_{n,M}f}_\infty
\lesssim
\left(
\frac{\log M}{\log\log M}
\right)^{-\alpha/d}.
}
\]

For a corrected Taylor approximation with
\[
a=\frac{r+1}{2},
\]
one obtains
\[
\boxed{
\norm{f-P_{n,M}f}_\infty
\lesssim
\left(
\frac{\log M}{\log\log M}
\right)^{-(r+1)/2}.
}
\]

When the hypotheses of Section~\ref{sec:raw-particle} or
Section~\ref{sec:root-n-taylor} genuinely give the uncorrected exponent
\(a=1\), one has
\[
\boxed{
\norm{f-P_{n,M}f}_\infty
\lesssim
\frac{\log\log M}{\log M}.
}
\]

\section{Conclusion}

There are three different sources of approximation error.

\begin{enumerate}[label=\arabic*.]
\item
The empirical-measure error:
\[
W_q(\widehat\mu_n,\mu),
\]
which makes a metric Hölder remainder or a \(C_W^\alpha\) Taylor
remainder contribute the high-dimensional rate
\[
n^{-\alpha/d}.
\]
For the raw operator this term must be combined with the weak bias, as
in Section~\ref{sec:raw-particle}.

\item
The weak particle bias:
\[
\frac{A_1f}{n}+\frac{A_2f}{n^2}+\cdots,
\]
which causes a generic uncorrected approximation to saturate at
\(1/n\) once a weak expansion makes the remaining error
\(o(n^{-1})\). Higher-order corrections remove these contractions and
are required to exploit additional Taylor regularity beyond this
threshold.

\item
The approximation of the kernel \(\Phi_n\), which is a function of
\(nd\) scalar variables. Under the full tensor convention, the number
of coefficients grows exponentially with \(nd\) at fixed
per-coordinate degree; other conventions still face a deteriorating
dimension-dependent approximation exponent.
\end{enumerate}

As a function of the interaction order \(n\), one can obtain algebraic
rates such as
\[
n^{-\alpha/d}
\qquad\text{or}\qquad
n^{-(r+1)/2}.
\]
Under the generic kernel-approximation model of
Section~\ref{sec:two-level}, the construction gives the logarithmic
parameter-count upper bound
\[
\left(
\frac{\log M}{\log\log M}
\right)^{-a}.
\]

To obtain an algebraic rate in \(M\) by this route, one needs additional
kernel structure, such as symmetry reduction, sparse interactions,
dependence on finitely many moments, low-rank tensor decompositions, or
kernels that are already finite-order polynomials, such as squared MMD
and pair-interaction energies.

\begin{thebibliography}{99}

\bibitem{babu}
G.~J. Babu,
\emph{A note on Dubins' theorem},
Canadian Mathematical Bulletin \textbf{33}(4) (1990), 416--418.

\bibitem{bertucci}
C.~Bertucci,
\emph{Analysis on spaces of measures},
arXiv:2607.19939, 2026.

\bibitem{carmona-delarue}
R.~Carmona and F.~Delarue,
\emph{Probabilistic Theory of Mean Field Games with Applications I:
Mean Field FBSDEs, Control, and Games},
Probability Theory and Stochastic Modelling, vol.~83, Springer, 2018.

\bibitem{chassagneux-szpruch-tse}
J.-F.~Chassagneux, {\L}.~Szpruch, and A.~Tse,
\emph{Weak quantitative propagation of chaos via differential calculus
on the space of measures},
The Annals of Applied Probability \textbf{32}(3) (2022), 1929--1969.

\bibitem{dereich-scheutzow-schottstedt}
S.~Dereich, M.~Scheutzow, and R.~Schottstedt,
\emph{Constructive quantization: approximation by empirical measures},
Annales de l'Institut Henri Poincaré, Probabilités et Statistiques
\textbf{49}(4) (2013), 1183--1203.

\bibitem{fournier-guillin}
N.~Fournier and A.~Guillin,
\emph{On the rate of convergence in Wasserstein distance of the
empirical measure},
Probability Theory and Related Fields \textbf{162} (2015), 707--738.

\bibitem{hoeffding}
W.~Hoeffding,
\emph{A class of statistics with asymptotically normal distribution},
The Annals of Mathematical Statistics \textbf{19}(3) (1948), 293--325.

\bibitem{richardson}
L.~F. Richardson,
\emph{The approximate arithmetical solution by finite differences of
physical problems involving differential equations, with an application
to the stresses in a masonry dam},
Philosophical Transactions of the Royal Society of London, Series A
\textbf{210} (1911), 307--357.

\bibitem{remez}
E.~Ya. Remez,
\emph{Sur la détermination des polynômes d'approximation de degré
donnée},
Communications de la Société Mathématique de Kharkov \textbf{10} (1934),
41--63.

\bibitem{schnabl}
R.~Schnabl,
\emph{Zur Approximation durch Bernsteinpolynome auf gewissen Räumen von
Wahrscheinlichkeitsmaßen},
Mathematische Annalen \textbf{180} (1969), 326--330.

\bibitem{serfling}
R.~J. Serfling,
\emph{Approximation Theorems of Mathematical Statistics},
Wiley, 1980.

\bibitem{von-mises}
R.~von Mises,
\emph{On the asymptotic distribution of differentiable statistical
functions},
The Annals of Mathematical Statistics \textbf{18}(3) (1947), 309--348.

\end{thebibliography}

\end{document}
